\par\Needspace{4\baselineskip}
\begin{center}
\hypertarget{native-059-article041}{}\pdfbookmark[2]{41.}{article-041}41.
\end{center}
Es bleibt noch, auch für die umgekehrte Aufgabe, nämlich die Bestimmung der Zeit aus der wahren Anomalie, einen einfacheren Rechenweg anzugeben. Dazu haben wir unserer Tafel eine neue Spalte für $T$ hinzugefügt. Zunächst berechnet man $T$ aus $v$ nach der Formel $T=\frac{1-e}{1+e}\operatorname{tang}\tfrac{1}{2}v^2$. Mit dem Argument $T$ entnimmt man unserer Tafel $A$ und $\log B$ oder, was genauer und sogar bequemer ist, $C$ und $\log B$, und erhält daraus $A$ nach der Formel $A=\frac{(1+C)T}{1+\frac{4}{5}T}$. Schließlich erhält man aus $A$ und $B$ den Wert $t$ nach Formel 1) aus Art. 37. Will man auch hier die \textsc{Barker}sche Tafel verwenden, so erleichtert sie bei dieser umgekehrten Aufgabe die Rechnung weniger. Man erhält dann unabhängig von $A$ unmittelbar
\[
\operatorname{tang}\tfrac{1}{2}w=\operatorname{tang}\tfrac{1}{2}v\,\mathord{.}\,\BareRoot{\frac{1+C}{\gamma\gamma(1+\frac{4}{5}T)}},
\]
und daraus die Zeit $t$ durch Multiplikation der mittleren Bewegung, die in der \textsc{Barker}schen Tafel der wahren Anomalie $w$ entspricht, mit $\frac{B}{\alpha}\cdot$
